% source: https://github.com/pfradin/luadraw/discussions/307#discussioncomment-18047899 (pfradin, block 1)
\documentclass[border=2 mm]{standalone}
\usepackage[svgnames]{xcolor}
\usepackage[3d]{luadraw}
\usepackage{fourier-otf}
\begin{document}
%%https://tex.stackexchange.com/questions/739936/fill-area-between-hyperbola-and-line-with-pgfplots
\begin{luadraw}{name=rotline}
local ld = luadraw
local pt3d = ld.pt3d
local cpx = ld.cpx
local sqrt = math.sqrt
local O, I, J, K = pt3d.Origin, pt3d.vecI, pt3d.vecJ, pt3d.vecK

local g = ld.graph3d:new{ window3d={-10,10,-10,10,-10,10}, window={-11,11,-11,11},
    viewdir={15,70}, size={12,12}}  
local f = function(x) return sqrt(3)/4*sqrt(256-x^2)  end
local h = function(x) return sqrt(100-x^2) end
local L = ld.domain3(f,h,-10,10) -- see documentation
L = ld.map(pt3d.toPoint3d,L) -- conversion 2D -> 3D
local S = ld.rotline(L,{O,I},-270,0, {close=true, nbdots=35})

g:Dfacet(S, {color="Tomato"})
-- close the two ends
g:Dfacet({L, ld.rotate3d(L,-270,{O,I})} , {color="Tomato", twoside=false})
g:Show()
\end{luadraw}
\end{document} 
